\documentclass[../../main.tex]{subfiles}
\begin{document}

\begin{flushright}
\footnotesize\textit{【自動文字起こし・要確認】}
\end{flushright}

{\bf [解]}

$(m,n)\neq (0,0)$の時，与えられた関数は
\begin{align}
f(m,n)=\dfrac{1}{2}(m+n)(m+n+1)+n \label{eq:1}
\end{align}
である．$f(0,0)=0$は\cref{eq:1}でも成立するので，結局$m,n\in\mathbb{Z}_{\ge0}$に対して\cref{eq:1}で良い．


\paragraph{(1)}

題意の条件$f(m,n) \le 5$に\cref{eq:1}を代入すると
\begin{align}
\dfrac{1}{2}(m+n)(m+n+1)+n&\le5 \\
\therefore 
(m+n)(m+n+1)+2n&\le 10 \quad\cdots\text{①} \label{eq:2}
\end{align}
である．$t=m+n$とすると$n\ge 0$だから$t(t+1)\le10$が必要である．
これを満たす$0 \le t\in\mathbb{N}$は$t=0,1,2$の三通りだから以下それぞれの場合について考える．

\subparagraph{$0^{\circ}$ $t=0$}

$t=0$となるのは$(m,n)=(0,0)$の時のみで，これは\cref{eq:2}を満たして十分である．

\subparagraph{$1^\circ$ $t=1$}

$t=1$となるのは$(m,n)=(0,1),(1,0)$の時のみで，これは\cref{eq:2}を満たして十分である．

\subparagraph{$2^\circ$ $t=2$}

$t=2$となるのは$(m,n)=(0,2),(1,1),(2,0)$の時のみで，これは\cref{eq:2}を満たして十分である．
\casesend

以上から，もとめる$(m,n)$は
\begin{align}
(m,n)=(0,0),(0,1),(0,2),(1,0),(1,1),(2,0)
\end{align}
で，$mn$平面に図示すると下図黒丸となる．

\begin{figure}[htb]
\begin{center}
\begin{tikzpicture}[scale=0.9, >=stealth]
  % 座標軸
  \draw[->] (-0.3,0) -- (3,0) node[right] {$m$};
  \draw[->] (0,-0.3) -- (0,3) node[above] {$n$};
  % f(m,n)<=5 をみたす点 (m,n) (t=m+nが1または2の場合)
  \foreach \p in {(0,0),(0,1),(0,2),(1,0),(1,1),(2,0)}{
    \fill \p circle (1.33pt);
  }
  \node[below] at (1,0) {$1$};
  \node[below] at (2,0) {$2$};
  \node[left] at (0,1) {$1$};
  \node[left] at (0,2) {$2$};
\end{tikzpicture}
\end{center}
\caption{条件を満たす$(m,n)$の図示（黒丸）}
\end{figure}

\paragraph{(2)}

$f(m,n)=f(m',n')$の場合を考える．\cref{eq:1}を代入すると
\begin{align}
\dfrac{1}{2}(m+n)(m+n+1)+n&=\dfrac{1}{2}(m'+n')(m'+n'+1)+n' \label{eq:3}
\end{align}
だから，$t=m+n, t'=m'+n'$とおくと
\begin{align}
t(t+1)+2n&=t'(t'+1)+2n' \label{eq:4}
\end{align}
である．

ここで，$g(x,y)=x^2+x+2y\ (0\le y\le x)$とおくと，これは$y$の単調増加関数だから
\begin{align}
x^2+x\le g(x,y)\le x^2+3x \quad\cdots\text{③}
\end{align}
となる．この両辺は$0\le x$では$x$についての単調増加関数である．

$x=t$として，
\begin{align}
  t^2+t\le f(t,n)\le t^2+3t 
\end{align}
であり，$\{(t+1)^2+(t+1)\}-(t^2+3t)=2>0$より，$0 \le n' \le t+1$なる自然数$n'$に対して
\begin{align}
 & t^2+t\le f(t,n)\le t^2+3t < \{(t+1)^2+(t+1)\} \le f(t+1,n') \\
\therefore
 & f(t, n) < f(t+1,n')
\end{align}
だから，異なる2自然数$t,t'$に対して$g(t,n)=g(t',n')$となることはない．
$f(t,n)=g(t,n)$だから\cref{eq:3}が成立するには$t=t'$が必要である．
\cref{eq:4}よりこの時$n=n'$とすれば十分であり，$m=m'$も成立する．

以上から，$f(m,n)=f(m',n')$ならば，$(m,n)=(m',n')$となるので題意は示された．


\newpage

\end{document}
